\documentclass[10pt,a4paper]{amsart}
\usepackage{amsmath,amssymb,amsthm,graphicx}
\usepackage{verbatim}
\newenvironment{reference}{\comment}{\endcomment}
%\newenvironment{reference}{}{}
\newenvironment{slogan}{\comment}{\endcomment}
\newenvironment{history}{\comment}{\endcomment}

% For commutative diagrams we use Xy-pic
\usepackage[all]{xy}

% We use 2cell for 2-commutative diagrams.
\xyoption{2cell}
% Theorem environments.
%
\theoremstyle{plain}
\newtheorem{theorem}[subsection]{Theorem}
\newtheorem{proposition}[subsection]{Proposition}
\newtheorem{lemma}[subsection]{Lemma}

\theoremstyle{definition}
\newtheorem{definition}[subsection]{Definition}
\newtheorem{example}[subsection]{Example}
\newtheorem{exercise}[subsection]{Exercise}
\newtheorem{situation}[subsection]{Situation}

\theoremstyle{remark}
\newtheorem{remark}[subsection]{Remark}
\newtheorem{remarks}[subsection]{Remarks}

\numberwithin{equation}{subsection}

% Macros
%
\def\lim{\mathop{\mathrm{lim}}\nolimits}
\def\colim{\mathop{\mathrm{colim}}\nolimits}
\def\Spec{\mathop{\mathrm{Spec}}}
\def\Hom{\mathop{\mathrm{Hom}}\nolimits}
\def\Ext{\mathop{\mathrm{Ext}}\nolimits}
\def\SheafHom{\mathop{\mathcal{H}\!\mathit{om}}\nolimits}
\def\SheafExt{\mathop{\mathcal{E}\!\mathit{xt}}\nolimits}
\def\Sch{\mathit{Sch}}
\def\Mor{\mathop{\mathrm{Mor}}\nolimits}
\def\Ob{\mathop{\mathrm{Ob}}\nolimits}
\def\Sh{\mathop{\mathit{Sh}}\nolimits}
\def\NL{\mathop{N\!L}\nolimits}
\def\CH{\mathop{\mathrm{CH}}\nolimits}
\def\proetale{{pro\text{-}\acute{e}tale}}
\def\etale{{\acute{e}tale}}
\def\QCoh{\mathit{QCoh}}
\def\Ker{\mathop{\mathrm{Ker}}}
\def\Im{\mathop{\mathrm{Im}}}
\def\Coker{\mathop{\mathrm{Coker}}}
\def\Coim{\mathop{\mathrm{Coim}}}

% Boxtimes
%
\DeclareMathSymbol{\boxtimes}{\mathbin}{AMSa}{"02}

%
% Macros for moduli stacks/spaces
%
\def\QCohstack{\mathcal{QC}\!\mathit{oh}}
\def\Cohstack{\mathcal{C}\!\mathit{oh}}
\def\Spacesstack{\mathcal{S}\!\mathit{paces}}
\def\Quotfunctor{\mathrm{Quot}}
\def\Hilbfunctor{\mathrm{Hilb}}
\def\Curvesstack{\mathcal{C}\!\mathit{urves}}
\def\Polarizedstack{\mathcal{P}\!\mathit{olarized}}
\def\Complexesstack{\mathcal{C}\!\mathit{omplexes}}
% \Pic is the operator that assigns to X its picard group, usage \Pic(X)
% \Picardstack_{X/B} denotes the Picard stack of X over B
% \Picardfunctor_{X/B} denotes the Picard functor of X over B
\def\Pic{\mathop{\mathrm{Pic}}\nolimits}
\def\Picardstack{\mathcal{P}\!\mathit{ic}}
\def\Picardfunctor{\mathrm{Pic}}
\def\Deformationcategory{\mathcal{D}\!\mathit{ef}}
\usepackage{cleveref}
\crefname{originallemma}{original lemma}{original lemmas}
\Crefname{originallemma}{Original lemma}{Original lemmas}
\makeatletter
\newlabel{algebra-lemma-K0-local}{{00JJ}{source}{}{sourcepointer-1}{}}
\newlabel{algebra-lemma-K0-local@cref}{{[originallemma][0][]00JJ}{[1][0][]source}{}{}{}}
\makeatother
\crefname{originalreference}{original reference}{original references}
\Crefname{originalreference}{Original reference}{Original references}
\makeatletter
\newlabel{lemma-affine-cap-s-open}{{01PV}{source}{}{sourcepointer-2}{}}
\newlabel{lemma-affine-cap-s-open@cref}{{[originalreference][0][]01PV}{[1][0][]source}{}{}{}}
\makeatother
\crefname{originalreference}{original reference}{original references}
\Crefname{originalreference}{Original reference}{Original references}
\makeatletter
\newlabel{lemma-locally-free-module}{{05JM}{source}{}{sourcepointer-3}{}}
\newlabel{lemma-locally-free-module@cref}{{[originalreference][0][]05JM}{[1][0][]source}{}{}{}}
\makeatother
\crefname{originallemma}{original lemma}{original lemmas}
\Crefname{originallemma}{Original lemma}{Original lemmas}
\makeatletter
\newlabel{schemes-lemma-affine-separated}{{01KN}{source}{}{sourcepointer-4}{}}
\newlabel{schemes-lemma-affine-separated@cref}{{[originallemma][0][]01KN}{[1][0][]source}{}{}{}}
\makeatother
\crefname{originallemma}{original lemma}{original lemmas}
\Crefname{originallemma}{Original lemma}{Original lemmas}
\makeatletter
\newlabel{schemes-lemma-characterize-quasi-separated}{{01KO}{source}{}{sourcepointer-5}{}}
\newlabel{schemes-lemma-characterize-quasi-separated@cref}{{[originallemma][0][]01KO}{[1][0][]source}{}{}{}}
\makeatother
\crefname{originallemma}{original lemma}{original lemmas}
\Crefname{originallemma}{Original lemma}{Original lemmas}
\makeatletter
\newlabel{schemes-lemma-quasi-coherent-affine}{{01IA}{source}{}{sourcepointer-6}{}}
\newlabel{schemes-lemma-quasi-coherent-affine@cref}{{[originallemma][0][]01IA}{[1][0][]source}{}{}{}}
\makeatother
\crefname{originallemma}{original lemma}{original lemmas}
\Crefname{originallemma}{Original lemma}{Original lemmas}
\makeatletter
\newlabel{schemes-lemma-valuative-criterion-separatedness}{{01L0}{source}{}{sourcepointer-7}{}}
\newlabel{schemes-lemma-valuative-criterion-separatedness@cref}{{[originallemma][0][]01L0}{[1][0][]source}{}{}{}}
\makeatother
\title{A scheme covered by affine opens of the form X\_s for sections of powers of an invertible sheaf is separated}
\author{The Stacks Project Authors}
\date{stacks-project:properties.tex:local-snapshot-20261012; GNU FDL 1.2 or later}
\begin{document}
\maketitle
\noindent\textit{Editorial note.} All delivered mathematical text is the author's own, copied from the pinned Stacks Project chapter file properties.tex; no statement or proof step was rewritten, reordered or removed.\par\medskip
\noindent\textit{Editorial note.} The counted proof is the author's complete argument: the derivation of quasi-separatedness from the affineness of the intersections, the setup of the valuative criterion with the passage to a common invertible sheaf, and the comparison of the free rank one modules inside their base changes to the fraction field concluding that both images lie in one affine open.\par\medskip
\noindent\textit{Editorial note.} The label lemma-affine-s-opens-cover-quasi-separated names the counted statement, and the proof environment immediately following it in the source is delivered verbatim as its proof.\par\medskip
\noindent\textit{Editorial note.} The counted statement is carried by the author's own lemma environment in the chapter source.\par\medskip
\noindent\textit{Editorial note.} Pre-open gate: python operations/agent-authoring/tools/precard\_gate.py --source properties.tex --label lemma-affine-s-opens-cover-quasi-separated --start 4579 --end 4640 --local-id 28040 returned RESULT=OK (L28040 free, occupancy FREE, public interval clear). Candidate taken from sources/stacks-properties/candidates-halfA.json.\par\medskip
\noindent\textit{Editorial note.} The statement and proof are copied from the pinned chapter file; the cited inputs are kept as external pointers with their canonical Stacks tags.\par\medskip
\noindent\textit{Allowed standard prerequisites (editorial).}\begin{itemize}
\item Intersections of two opens of the form X\_s are affine. Same chapter, the lemma that the intersection of the nonvanishing loci of two sections of powers of an invertible sheaf is affine. Supplies the affineness of the pairwise intersections, hence quasi-separatedness.
\item Valuative criterion for separatedness. Schemes, the lemma characterizing separated morphisms through the valuative criterion. Supplies the reduction of separatedness to the comparison of two morphisms from the spectrum of a valuation ring agreeing over the fraction field.
\item Locally free rank one modules over a local ring are free. Same chapter, the lemma identifying the pullback of an invertible sheaf along a morphism from an affine scheme with a locally free rank one module, and the fact that such a module over a local ring is free. Supply the identification of the two pullbacks with free rank one modules over the valuation ring, hence with submodules of their base change to the fraction field.
\item Affine schemes are separated. Schemes, the lemma that affine schemes are separated. Supplies the final step concluding the two morphisms agree on the affine open.
\end{itemize}\medskip
\ifdefined\hypertarget\hypertarget{sourcepointer-1}{}\fi
\noindent\textit{Source reference.} Original lemma 00JJ: sources/stacks-chow/tags.txt, source lines 717--717. Prerequisite pointer invoked at proof line(s) 4619; it names a standard fact of a sibling chapter that the counted argument uses, and every step of the counted proof built on it is retained verbatim in the delivered body.\par\medskip
\ifdefined\hypertarget\hypertarget{sourcepointer-2}{}\fi
\noindent\textit{Source reference.} Original reference 01PV: sources/stacks-properties/pinned-original-source.tex, source lines 4505--4505. Prerequisite pointer invoked at proof line(s) 4590; it names a standard fact of the same chapter that the counted argument uses, and every step of the counted proof built on it is retained verbatim in the delivered body.\par\medskip
\ifdefined\hypertarget\hypertarget{sourcepointer-3}{}\fi
\noindent\textit{Source reference.} Original reference 05JM: sources/stacks-properties/pinned-original-source.tex, source lines 2852--2852. Prerequisite pointer invoked at proof line(s) 4617; it names a standard fact of the same chapter that the counted argument uses, and every step of the counted proof built on it is retained verbatim in the delivered body.\par\medskip
\ifdefined\hypertarget\hypertarget{sourcepointer-4}{}\fi
\noindent\textit{Source reference.} Original lemma 01KN: sources/stacks-chow/tags.txt, source lines 2056--2056. Prerequisite pointer invoked at proof line(s) 4638; it names a standard fact of a sibling chapter that the counted argument uses, and every step of the counted proof built on it is retained verbatim in the delivered body.\par\medskip
\ifdefined\hypertarget\hypertarget{sourcepointer-5}{}\fi
\noindent\textit{Source reference.} Original lemma 01KO: sources/stacks-chow/tags.txt, source lines 2057--2057. Prerequisite pointer invoked at proof line(s) 4592; it names a standard fact of a sibling chapter that the counted argument uses, and every step of the counted proof built on it is retained verbatim in the delivered body.\par\medskip
\ifdefined\hypertarget\hypertarget{sourcepointer-6}{}\fi
\noindent\textit{Source reference.} Original lemma 01IA: sources/stacks-chow/tags.txt, source lines 1971--1971. Prerequisite pointer invoked at proof line(s) 4615; it names a standard fact of a sibling chapter that the counted argument uses, and every step of the counted proof built on it is retained verbatim in the delivered body.\par\medskip
\ifdefined\hypertarget\hypertarget{sourcepointer-7}{}\fi
\noindent\textit{Source reference.} Original lemma 01L0: sources/stacks-chow/tags.txt, source lines 2069--2069. Prerequisite pointer invoked at proof line(s) 4598; it names a standard fact of a sibling chapter that the counted argument uses, and every step of the counted proof built on it is retained verbatim in the delivered body.\par\medskip
\begin{lemma}
\label{lemma-affine-s-opens-cover-quasi-separated}
Let $X$ be a scheme and $\mathcal{L}$ be an invertible $\mathcal{O}_X$-module.
Assume for every point $x$ of $X$ there exists $n \geq 1$ and
$s \in \Gamma(X, \mathcal{L}^{\otimes n})$ such that
$x \in X_s$ and $X_s$ is affine. Then $X$ is separated.
\end{lemma}

\begin{proof}
We show first that $X$ is quasi-separated.  By assumption we
can find a covering of $X$ by affine opens of the form
$X_s$.  By Lemma \ref{lemma-affine-cap-s-open}, the
intersection of any two such sets is affine, so Schemes, Lemma
\ref{schemes-lemma-characterize-quasi-separated} implies
that $X$ is quasi-separated.

\medskip\noindent
To show that $X$ is separated, we can use the valuative
criterion, Schemes, Lemma
\ref{schemes-lemma-valuative-criterion-separatedness}.
Thus, let $A$ be a valuation ring with fraction field $K$
and consider two morphisms $f, g : \Spec(A) \to X$ such that
the two compositions $\Spec(K) \to \Spec(A) \to X$ agree.
As $A$ is local, there exists $p, q \ge 1$, $s \in \Gamma(X,
\mathcal{L}^{\otimes p})$, and $t \in \Gamma(X,
\mathcal{L}^{\otimes q})$ such that $X_s$ and $X_t$ are
affine, $f(\Spec A) \subseteq X_s$, and $g(\Spec A)
\subseteq X_t$.  We now replace $s$ by $s^q$, $t$ by $t^p$,
and $\mathcal{L}$ by $\mathcal{L}^{\otimes pq}$.  This is
harmless as $X_s = X_{s^q}$ and $X_t = X_{t^p}$, and now $s$
and $t$ are both sections of the same sheaf $\mathcal{L}$.

\medskip\noindent
The quasi-coherent module $f^*\mathcal{L}$ corresponds to an $A$-module $M$ and
$g^*\mathcal{L}$ corresponds to an $A$-module $N$ by our
classification of quasi-coherent modules over affine schemes
(Schemes, Lemma \ref{schemes-lemma-quasi-coherent-affine}).
The $A$-modules $M$ and $N$ are locally free of rank
$1$ (Lemma \ref{lemma-locally-free-module}) and as $A$ is
local they are free (Algebra, Lemma
\ref{algebra-lemma-K0-local}).  Therefore we may identify
$M$ and $N$ with $A$-submodules of $M \otimes_A K$ and $N
\otimes_A K$.  The equality $f|_{\Spec(K)} = g|_{\Spec(K)}$
determines an isomorphism $\phi \colon M \otimes_A K \to N
\otimes_A K$.

\medskip\noindent
Let $x \in M$ and $y \in N$ be the elements corresponding to
the pullback of $s$ along $f$ and $g$, respectively.  These
satisfy $\phi(x \otimes 1) = y \otimes 1$.  The image of $f$
is contained in $X_s$, so $x \not\in \mathfrak{m}_A M$, that
is, $x$ generates $M$.  Hence $\phi$ determines an
isomorphism of $M$ with the submodule of $N$ generated by
$y$.  Arguing symmetrically using $t$, $\phi^{-1}$
determines an isomorphism of $N$ with a submodule of $M$.
Consequently $\phi$ restricts to an isomorphism of $M$ and
$N$.  Since $x$ generates $M$, its image $y$ generates
$N$, implying $y \not\in \mathfrak{m}_A N$.  Therefore
$g(\Spec(A)) \subseteq X_s$.  Because $X_s$ is affine, it is
separated by Schemes, Lemma \ref{schemes-lemma-affine-separated},
and we conclude $f = g$.
\end{proof}

\end{document}
